\section{Two direct constructions for small beta sets}\label{sec:small}

\subsection{Saving zeros by allowing more than one value}

To prove that at least one of $\beta(4),\beta(6)$ is irrational, it
would suffice to construct nonzero integer linear forms in these two
values tending to zero. Under joint rationality, a single fixed
denominator would make all these forms integers. There is no need to
cancel either retained coefficient. The same applies to the three-value
set $\{\beta(4),\beta(6),\beta(8)\}$.

The differentiated rational-function method, as used by
Krisman~\cite[Sections~2--4, especially Lemma~4.6]{Krisman}, provides
a natural test. Differentiating twice shifts the first even period
from $\beta(2)$ to $\beta(4)$. Its usual numerator zeros have
multiplicity at least three, so they survive the differentiation.
Could simple zeros save enough numerator degree to compensate for
the rational boundary terms they leave? This is the question tested
in \rec{R29}, following \rec{R28}.

Let $n,r,v$ be positive odd integers, and let $A>rv$ be even. Define
\begin{equation}\label{eq:small-rational}
 R_n(t)=(2^nn!)^{A-rv}
 \frac{\displaystyle\prod_{j=0}^{rn-1}
               (t-(r-1)n+2j+1)^v}
      {\displaystyle\prod_{h=0}^n(t+2h)^A}.
\end{equation}
The factorial normalization cancels the $n\log n$ terms in Stirling
estimates. It does not change the primitive integer form, except
possibly its sign. Reflection gives $R_n(-2n-t)=-R_n(t)$. If
\[
 R_n(t)=\sum_{h=0}^n\sum_{i=1}^A\frac{a_{i,h}}{(t+2h)^i},
\]
then $a_{i,n-h}=(-1)^{i+1}a_{i,h}$ and
$\sum_h(-1)^ha_{i,h}=0$ for odd $i$.

Put $u=(r-1)n/2$ and $S_K(s)=\sum_{j=0}^{K-1}(-1)^j/(2j+1)^s$.
The full sum and its tail are
\[
 F_n=\sum_{j\ge0}(-1)^jR_n''(2j+1)/2,\qquad
 T_n=\sum_{j\ge u}(-1)^jR_n''(2j+1)/2.
\]
Partial fractions give
\begin{align}\label{eq:small-form}
 T_n&=q_u+\sum_{s=4,6,\ldots,A+2}\rho_s\beta(s),\nonumber\\
 \rho_{i+2}&=\binom{i+1}2\sum_h(-1)^ha_{i,h},\qquad
 q_u=-\sum_{h,i}(-1)^h\binom{i+1}2a_{i,h}S_{u+h}(i+2).
\end{align}
The constant sum includes \emph{all} $i$, including odd $i$. Their
aggregate beta coefficients vanish, but their rational finite sums
need not. Setting $u=0$ gives the full sum and $q_0$, so
\begin{equation}\label{eq:small-prefix}
 q_0-q_u=\sum_{j=0}^{u-1}(-1)^jR_n''(2j+1)/2.
\end{equation}
For $v\ge3$ this prefix is zero. For $v=1$ it generally is not.

\subsection{The zeros remove large denominator primes by grouping}

The zeros serve an arithmetic purpose as well as making terms vanish.
Suppose a proper rational function
$V(t)=\sum_{h,i}c_{i,h}(t+2h)^{-i}$ has $p$-integral coefficients
and vanishes at each negative odd integer from $-2n+1$ to $-1$,
where $p>n$ is odd. In the rational constant of its alternating sum,
group terms by that negative odd integer $x$. The group is
\[
 -\chi(x)\sum_{h:x+2h>0}\sum_i\frac{c_{i,h}}{(x+2h)^i},
 \qquad \chi(x)=(-1)^{(x-1)/2}.
\]
The sum over all $h$ is $V(x)=0$, so the complementary range
$x+2h<0$ represents the same group with the opposite sign.
A denominator divisible by $p$ cannot occur in both ranges, since
that would give $p\mid2(h-h')$ with $0<|h-h'|\le n<p$.
At least one representation is therefore $p$-integral.

For \eqref{eq:small-rational}, the coefficients are $p$-integral for
odd $p>n$: local expansion only divides by differences $2(k-h)$.
The required zeros survive for $V=R_n''/2$ when $v\ge3$, and for
$V=WR_n$ without differentiation when $W$ is an integer polynomial
and the product is proper. These constructions have no odd denominator
primes above $n$ in their rational constants. The simple-zero derivative
does not enjoy this cancellation.

\subsection{An exact residue formula for the extra cost}

Take $v=1$ and an odd prime satisfying
\[
 \max(rn,A+2)<p\le(r+1)n-1,\qquad
 t=\frac{(r+1)n-1-p}{2}.
\]
In the finite sums for $q_u$, the only positive denominator divisible
by $p$ is $p$ itself. It occurs at poles $h=n-t+j$, $0\le j\le t$.
At such a pole exactly one numerator factor is $-p$. If $K_h$ is
the product of the other factors divided by the other pole factors,
local expansion gives
$a_{A,h}/p\equiv-K_h$ and $a_{A-1,h}\equiv K_h\pmod p$.
After multiplying $q_u$ by $p^{A+1}$, all lower orders vanish.
The two surviving binomial coefficients differ by $A$, and direct
factorial evaluation gives
\begin{equation}\label{eq:tail-residue}
 p^{A+1}q_u\equiv
 \frac A2\chi(p)(-1)^{n-t}(n!)^{A-r}
 \sum_{j=0}^t\frac{j!(rn-1-j)!}
 {((n-t+j)!)^A((t-j)!)^A}\pmod p.
\end{equation}
Every factorial here has argument below $p$. This proves
$v_p(q_u)\ge-(A+1)$, with equality whenever the residue is nonzero.
Since all beta coefficients are $p$-integral, equality forces
$p^{A+1}$ into the full primitive multiplier.

For $r\ge3$, the full constant $q_0$ is $p$-integral because its
finite denominators are at most $2n-1<p$. Thus these primes are a
new cost of removing the prefix. For example, at $n=25,A=4,r=3$,
each of $79,83,89,97$ occurs to exponent five in the primitive tail
multiplier and is absent from the full-sum cost. The nonvanishing
qualification in \eqref{eq:tail-residue} matters: among the $67$
recorded comparisons, one residue vanishes. At $n=9,A=2,r=1,p=11$,
the actual constant valuation is $-2$, rather than $-3$.

\subsection{Keep the zeros and cancel \texorpdfstring{$\beta(2)$}{beta(2)} instead}

A second test avoids differentiation. Form the full sums of $R_n$
and $(t+n)^2R_n$, with $v=1$. Both have periods through $\beta(A)$,
starting at $\beta(2)$, and their constants have the grouped
cancellation just proved. Multiplication at a pole gives
\[
 a^{(2)}_{i,h}=(n-2h)^2a_{i,h}
          +2(n-2h)a_{i+1,h}+a_{i+2,h}.
\]
Call the two complete forms $U_0,U_2$, and their $\beta(2)$
coefficients $b_0,b_2$. With $g_2=\cont(b_0,b_2)$, the combination
\[
 H_n=(b_2/g_2)U_0-(b_0/g_2)U_2
\]
uses primitive integer weights and cancels $\beta(2)$ exactly.
For $A=6$ it retains $\beta(4),\beta(6)$; for $A=8$ it retains the
three-value set. Its complete coefficient vector is then reduced to
primitive form once more. The weights depend on $n$, so the fixed
multiplier theorem of Section~\ref{sec:isolation} does not apply.

\subsection{The finite comparisons}

The differentiated tests use $n=1,3,5,9,17,25$. At $n=25$, the
natural logarithms of absolute values, divided by $n$, are as follows.
The primitive columns include the rational constant and the complete
coefficient gcd, as in Lemma~\ref{lem:constant-cost}.
\begin{center}
\begin{tabular}{@{}lrrrrrr@{}}
\toprule
Retained values&$A$&$r$&$v$&Raw tail&\shortstack{Primitive\\tail}&\shortstack{Primitive\\full sum}\\
\midrule
$\beta(4)$&2&1&1&$-.06$&7.48&7.48\\
$\beta(4),\beta(6)$&4&1&1&$-.13$&11.28&11.28\\
$\beta(4),\beta(6)$&4&3&1&$-2.85$&18.15&13.37\\
$\beta(4),\beta(6)$&4&1&3&$-.18$&8.96&8.96\\
$\beta(4),\beta(6),\beta(8)$&6&3&1&$-5.80$&22.79&17.35\\
$\beta(4),\beta(6),\beta(8)$&6&5&1&$-4.48$&37.78&20.36\\
$\beta(4),\beta(6),\beta(8)$&6&1&3&$-.26$&10.87&10.87\\
$\beta(4),\ldots,\beta(12)$&10&3&3&$-5.44$&15.69&15.69\\
\bottomrule
\end{tabular}
\end{center}
For $A=4,r=3$, removing the prefix saves about $3.81$ in raw size
but adds about $8.59$ in constant-denominator cost, in this normalization.
The resulting primitive form is larger. These are finite values, not
claimed limiting exponents.
The last row keeps five even values: with $r=v=3$, the degree condition
$A>rv$ requires $A\ge10$. It shows the larger support needed to retain
triple zeros on that longer numerator block.

The cancellation alternative gives the following primitive values.
\begin{center}
\begin{tabular}{lrrr}
\toprule
Retained support&$A$&$r$&$\log|\text{primitive }H_{25}|/25$\\
\midrule
$\beta(4)$&4&3&10.88\\
$\beta(4),\beta(6)$&6&1&12.39\\
$\beta(4),\beta(6)$&6&3&13.13\\
$\beta(4),\beta(6)$&6&5&20.53\\
$\beta(4),\beta(6),\beta(8)$&8&3&15.36\\
$\beta(4),\beta(6),\beta(8)$&8&5&22.03\\
$\beta(4),\beta(6),\beta(8)$&8&7&32.61\\
\bottomrule
\end{tabular}
\end{center}
For $A=6,r=3$, the two primitive input forms have rates about
$3.13$ and $3.10$. Eliminating $\beta(2)$ raises the result to $13.13$;
the primitive cancellation weights have up to $429$ bits. The $A=4$
control is the earlier adaptive-isolation family, up to scalar and sign.

All $115$ final primitive forms at the tested $n\ge3$ have absolute
value greater than one. Some first-index forms are small: the
$A=6,r=3,n=1$ cancellation gives
$731-182\beta(4)-552\beta(6)$. Its isolated smallness gives no
sequence tending to zero. The tests identify costs of these two
changes; they do not rule out a different small-set construction.
In particular, Krisman's shorter finite-set theorem uses a specifically
balanced rational function and its arithmetic savings. Testing the
simple families above does not test every choice in that method.

The checks reconstruct all partial fractions by a second algorithm,
retain every constant term, and compare all $222$ complete forms,
including intermediate ones. Numerical signs and sizes are enclosed
both through Arb Hurwitz zeta values and through an Euler-transformed
alternating series. The latter has positive differences
\[
 \Delta^jf_s(0)=\frac{4^j(j!)^2}{(2j+1)!}
 h_{s-1}\left(1,\tfrac13,\ldots,\tfrac1{2j+1}\right),
 \qquad f_s(k)=(2k+1)^{-s},
\]
where $h_d$ is the complete homogeneous symmetric polynomial.
Summing $2^{-j-1}\Delta^jf_s(0)$ and bounding its remaining positive
tail supplies an independent numerical enclosure.
